\documentclass[12pt,a4paper]{article}

% !TeX program = lualatex

% ------------------------------------------------
% Fonts and Unicode
% ------------------------------------------------
\usepackage{fontspec}
\setmainfont{Libertinus Serif}

% ------------------------------------------------
% Mathematics
% ------------------------------------------------
\usepackage{amsmath}
\usepackage{amssymb}

% ------------------------------------------------
% Layout
% ------------------------------------------------
\usepackage[
    a4paper,
    margin=1in
]{geometry}

% ------------------------------------------------
% Tables
% ------------------------------------------------
\usepackage{booktabs}
\usepackage{tabularx}

% ------------------------------------------------
% Quotes
% ------------------------------------------------
\usepackage{csquotes}

% ------------------------------------------------
% Links
% ------------------------------------------------
\usepackage[
    colorlinks=true,
    linkcolor=blue,
    citecolor=blue,
    urlcolor=blue
]{hyperref}

% ------------------------------------------------
% Other
% ------------------------------------------------
\usepackage{microtype}
\usepackage{enumitem}

% ------------------------------------------------
% Document information
% ------------------------------------------------

\title{
    The Solvability of Chess\\
    \large An Analysis of the Elon Musk--Chess.com Exchange
}

\author{Chat-gpt}
\date{September 2026}

% ------------------------------------------------
% Document
% ------------------------------------------------

\begin{document}

\maketitle

\begin{abstract}

In September 2026, Elon Musk and Chess.com engaged in a public
exchange on X concerning the computational complexity and eventual
solvability of chess. What initially appeared to be a brief social-media
argument developed into a discussion involving several distinct
mathematical concepts: the number of possible chess games, the number
of legal chess positions, computational tractability, and the meaning
of the word ``solved.''

This essay examines the exchange as both a technical disagreement and
a case study in how complex mathematical ideas are compressed into
short-form social-media arguments. Particular attention is given to
the distinction between the size of the chess game tree and the number
of reachable legal positions. The analysis also considers whether
``solving chess'' necessarily requires explicitly enumerating every
possible game or position.

\end{abstract}


\tableofcontents

\newpage


% ================================================================
\section{Introduction}
% ================================================================

Chess occupies an unusual position among games. Its rules are simple
enough to be explained in a few pages, yet the number of possible
positions and game sequences is so large that exhaustive analysis is
far beyond the practical capabilities of present-day computers.

Modern chess engines have nevertheless become extraordinarily strong.
Systems based on algorithms such as alpha--beta search, Monte Carlo
methods, evaluation functions, neural networks, and enormous
computational resources can already play chess at a level far beyond
ordinary human players.

This creates an interesting question:

\begin{quote}
    Is chess merely extremely difficult to solve, or is it ultimately
    solvable in a practical sense?
\end{quote}

The distinction became the subject of a public exchange between Elon
Musk and the Chess.com account on X in September 2026.

Musk argued that the number of genuinely useful moves in chess is
small and predicted that chess would eventually be ``fully solved''
in a manner analogous to checkers. Chess.com's response emphasized
the enormous number of possible games.

At first glance, the disagreement appears to concern a simple question
of arithmetic. In reality, the exchange mixes several different
questions about the structure of chess.

This essay separates those questions.


% ================================================================
\section{The Exchange}
% ================================================================

\subsection{Musk's Initial Claim}

On September 3, 2026, Musk argued on X that the number of moves in
chess that are not effectively useless is very small and that chess
would eventually be fully solved.

The important part of his argument was not merely that computers are
already stronger than humans. Rather, it was the extrapolation from
present computational superiority toward eventual mathematical
solvability.

That distinction matters.

A game can be computationally manageable for a machine in many
individual positions while remaining enormously difficult to solve
globally.


\subsection{Chess.com's Response}

Chess.com's official X account responded with the deliberately terse
phrase:

\begin{quote}
    ``Skill issue.''
\end{quote}

The exchange subsequently became more technical.

Chess.com emphasized the enormous number of possible chess games,
arguing that the game tree is vastly larger than the number of atoms
in the observable universe.

Musk objected to this comparison and redirected the discussion toward
the number of legal chess positions.

This was the point at which the argument became substantially more
interesting.


% ================================================================
\section{Two Different Numbers}
% ================================================================

The central mathematical confusion in the exchange comes from treating
two related quantities as though they were interchangeable.

They are not.

\subsection{Possible Games}

A chess game is a sequence of positions produced by successive legal
moves.

Even if the number of legal positions were relatively manageable,
the number of possible sequences through those positions could be
enormous.

A rough historical estimate commonly associated with Claude Shannon
places the number of possible chess games around

\[
10^{120}.
\]

This figure is often called the ``Shannon number.''

It should not be interpreted as an exact enumeration of every possible
game. It is an order-of-magnitude estimate illustrating the enormous
branching factor of chess.

\subsection{Legal Positions}

A different quantity is the number of legal positions that can occur
on a chessboard.

Estimates for the number of legal positions are dramatically smaller
than the estimated number of complete game sequences.

This distinction supports an important part of Musk's objection.

If the objective were literally to store every possible game sequence,
the task would be vastly larger than storing every position.

However, that does not mean that storing every legal position would
automatically solve chess.


% ================================================================
\section{What Does ``Solved'' Mean?}
% ================================================================

The word ``solved'' is doing a surprising amount of work in the debate.

There are several different meanings.

\subsection{Weakly Solved}

A game is weakly solved if the theoretical result from the initial
position is known under perfect play.

For example, if perfect play from the initial position necessarily
produces a draw, the game is weakly solved if that fact has been
mathematically established.

\subsection{Strongly Solved}

A game is strongly solved when the theoretically optimal result can
be determined from every reachable position.

This is a substantially stronger requirement.

For chess, strong solution would involve determining the optimal
outcome of arbitrary legal positions, not merely determining the
result from the initial board.

\subsection{Practical Solving}

There is also a third concept that is often implicit in discussions
about artificial intelligence:

\begin{quote}
    A system may be capable of playing essentially perfect chess
    without possessing a complete mathematical solution to chess.
\end{quote}

This distinction is crucial.

A sufficiently powerful engine could potentially become practically
unbeatable without humanity ever possessing a compact table containing
the theoretically optimal continuation from every possible position.


% ================================================================
\section{Game Trees and Compression}
% ================================================================

Musk's most interesting argument concerns the possibility that a future
artificial intelligence system could discover a much more efficient
representation of the problem.

Suppose, for illustration, that a game contains a set of states

\[
S = \{s_1,s_2,\ldots,s_n\}.
\]

A naive approach might attempt to examine every possible sequence of
moves.

But many different sequences can eventually reach the same position.
Consequently, analyzing every sequence independently would duplicate
work.

This is one reason that modern chess engines do not simply enumerate
every possible game.

They exploit the structure of the search space.

For example, if several different move sequences lead to the same
position \(s\), the value of \(s\) can potentially be reused:

\[
V(s) = f(s).
\]

Instead of calculating the consequences of every path independently,
an algorithm can calculate the value of the position and reuse that
result.

This is conceptually similar to dynamic programming.

Therefore:

\begin{quote}
    Number of possible games
    \neq
    computational work required to solve the game.
\end{quote}

The difference between those quantities is one of the most important
points obscured by the original social-media argument.


% ================================================================
\section{Why the Atom Comparison Is Both Useful and Limited}
% ================================================================

Chess.com's comparison between possible games and atoms in the
observable universe is rhetorically powerful because it gives a
human-scale intuition for an otherwise incomprehensible number.

However, it is not by itself an argument that chess cannot be solved.

Consider the following simplified analogy.

Suppose a problem has

\[
10^{120}
\]

possible paths but only

\[
10^{45}
\]

meaningfully distinct states.

The number of paths may be astronomically large while the underlying
state space is considerably smaller.

Furthermore, a mathematical solution does not necessarily require
explicitly listing every path.

A sufficiently compact theorem, algorithm, or representation could
in principle describe the behavior of an enormous number of states.

This is why the statement

\begin{quote}
    ``There are more possible games than atoms''
\end{quote}

should not be treated as equivalent to

\begin{quote}
    ``Chess cannot be solved.''
\end{quote}

The former concerns the size of a combinatorial space.

The latter is a claim about the existence and computational
tractability of a solution.


% ================================================================
\section{The Role of Artificial Intelligence}
% ================================================================

The development of increasingly powerful chess engines demonstrates
that enormous search spaces can be attacked without literally
enumerating every possibility.

Modern systems combine search, evaluation, learned representations,
and extensive prior computation.

This changes the practical question.

Instead of asking:

\begin{quote}
    ``Can a computer examine every possible game?''
\end{quote}

one can ask:

\begin{quote}
    ``Can a computer construct a sufficiently compressed representation
    of chess that determines optimal play?''
\end{quote}

Those are different problems.

The second question is much closer to the argument Musk was making.

Whether an artificial intelligence system could eventually discover
such a representation is an open technological and mathematical
question rather than something established by the number of legal
positions alone.


% ================================================================
\section{The Rhetoric of the X Exchange}
% ================================================================

The exchange is also interesting because the technical disagreement
was embedded inside the communication style of X.

Chess.com's initial ``skill issue'' response was intentionally
provocative rather than explanatory.

Musk responded by defending the underlying proposition and eventually
criticized the use of the observable-universe comparison.

The discussion subsequently became partly about the identity and
status of the person operating the Chess.com account. Musk referred
to the account operator as an ``intern,'' while the account operator
responded that they were a full-time employee.

This illustrates a recurring feature of technical arguments on social
media: the discussion can rapidly move between several levels.

\begin{enumerate}
    \item The original technical proposition.
    \item The evidence supporting that proposition.
    \item The credibility of the participants.
    \item The rhetorical framing of the disagreement.
    \item The audience's interpretation of the exchange.
\end{enumerate}

The result is entertaining, but it also makes careful technical
analysis more difficult.


% ================================================================
\section{Where the Arguments Intersect}
% ================================================================

The most interesting conclusion is that the two sides were not
necessarily answering exactly the same question.

\begin{table}[ht]
\centering
\begin{tabularx}{\textwidth}{lXX}
\toprule
Concept & What it describes & Relevance \\
\midrule
Game tree &
Possible sequences of moves &
Measures the enormous branching structure of chess \\

Legal positions &
Distinct reachable board states &
Measures the underlying state space \\

Perfect play &
Optimal decisions from positions &
Describes what a solution would need to determine \\

Practical engine strength &
Ability to choose strong moves &
Does not necessarily constitute a mathematical solution \\

Strong solution &
Optimal result from arbitrary legal positions &
A much stronger requirement than merely playing extremely well \\
\bottomrule
\end{tabularx}
\caption{Different concepts involved in the solvability question.}
\end{table}

Chess.com's argument demonstrates why chess is extraordinarily
difficult.

Musk's response highlights why the raw number of possible games is
not, by itself, sufficient to establish impossibility.

The deeper question lies between them:

\begin{quote}
    How much of the enormous chess state space can be compressed into
    a representation from which optimal play can be derived?
\end{quote}

That is a considerably more interesting problem than the original
social-media exchange suggests.


% ================================================================
\section{Conclusion}
% ================================================================

The Elon Musk--Chess.com exchange began as a short argument about
whether chess would eventually be solved, but it exposed a genuine
conceptual distinction.

The number of possible chess games and the number of legal chess
positions are not the same quantity. More importantly, neither number
alone determines whether chess can eventually be solved.

A complete solution could potentially exploit enormous amounts of
mathematical structure rather than explicitly enumerating every game.

At the same time, the existence of a much smaller set of legal
positions than possible game sequences does not make the problem
trivial. Determining optimal play over that state space remains an
extraordinarily difficult computational problem.

The most useful interpretation of the debate is therefore not that
one participant produced a single decisive number and the other did
not. Rather, the exchange accidentally exposed the distinction between
the size of a problem's search space and the amount of information
required to solve it.

That distinction extends well beyond chess.

\newpage

% ================================================================
\begin{thebibliography}{9}

\bibitem{complex}
Gomez, Jade.
\textit{Chess Website Humiliates Elon Musk on X: ``Skill Issue.''}
Complex, September 6, 2026.

\url{https://www.complex.com/pop-culture/a/jadegomez510/chess-com-elon-musk-x-beef-twitter}

\bibitem{usatoday}
Carry, Owen.
\textit{Elon Musk is feuding with Chess.com. And not for the first time.}
USA Today, September 4, 2026.

\url{https://www.usatoday.com/}

\bibitem{chessbharat}
Chess Bharat.
\textit{Will Chess Really Be Solved? Elon Musk vs. Chess.com Argument.}
September 9, 2026.

\url{https://chessbharat.org/elon-musk-vs-chess-com-will-chess-be-solved/}

\end{thebibliography}

\end{document}
